\chapter{Quantizing the Free Scalar Field}

\section{Classical modes}

The real Klein--Gordon equation is
\[
  (\p_t^2-\nabla^2+m^2)\phi=0.
\]
In a box of volume \(V=L^3\) with periodic boundary conditions, expand
\[
  \phi(t,\bm x)=\frac1{\sqrt V}\sum_{\bm k}q_{\bm k}(t)e^{\ii\bm k\cdot\bm x},
  \qquad
  \bm k=\frac{2\pi}{L}(n_x,n_y,n_z).
\]
Reality of \(\phi\) gives \(q_{-\bm k}=q_{\bm k}^*\).  Substituting into the equation of motion,
\[
  \ddot q_{\bm k}+\omega_{\bm k}^2q_{\bm k}=0,
  \qquad
  \omega_{\bm k}=\sqrt{|\bm k|^2+m^2}.
\]
Every Fourier mode is a harmonic oscillator.  This is the simplest bridge from ordinary quantum mechanics to QFT.

\section{Canonical commutation relations}

The conjugate momentum is
\[
  \pi(t,\bm x)=\p_t\phi(t,\bm x).
\]
Canonical quantization promotes \(\phi,\pi\) to operators and imposes
\[
  [\phi(t,\bm x),\pi(t,\bm y)]=\ii\delta^{(3)}(\bm x-\bm y),
\]
\[
  [\phi(t,\bm x),\phi(t,\bm y)]=0,\qquad
  [\pi(t,\bm x),\pi(t,\bm y)]=0.
\]
These are the field version of \([\hat x,\hat p]=\ii\).  The delta function appears because there is one oscillator at each point, or equivalently one oscillator for each mode.

\section{Creation and annihilation operators}

The field operator can be expanded as
\[
  \phi(x)
  =
  \int\frac{\dd^3p}{(2\pi)^3}
  \frac{1}{\sqrt{2E_{\bm p}}}
  \left[
  a_{\bm p}e^{-\ii p\cdot x}
  +
  a_{\bm p}^\dagger e^{\ii p\cdot x}
  \right],
\]
where
\[
  E_{\bm p}=\sqrt{|\bm p|^2+m^2},
  \qquad
  p\cdot x=E_{\bm p}t-\bm p\cdot\bm x.
\]
The conjugate momentum is
\[
  \pi(x)=\p_t\phi(x)
  =
  \int\frac{\dd^3p}{(2\pi)^3}
  (-\ii)\sqrt{\frac{E_{\bm p}}{2}}
  \left[
  a_{\bm p}e^{-\ii p\cdot x}
  -
  a_{\bm p}^\dagger e^{\ii p\cdot x}
  \right].
\]
The required commutator is obtained if
\[
  [a_{\bm p},a_{\bm q}^\dagger]=(2\pi)^3\delta^{(3)}(\bm p-\bm q),
\]
with all other \(a,a^\dagger\) commutators zero.

\derivation{
At equal time, compute the contribution of \(a_{\bm p}\) with \(a_{\bm q}^\dagger\) and the contribution of \(a_{\bm p}^\dagger\) with \(a_{\bm q}\).  The first gives
\[
  \int\frac{\dd^3p}{(2\pi)^3}
  \frac{-\ii}{2}
  e^{\ii\bm p\cdot(\bm x-\bm y)}.
\]
The second gives the same after changing \(\bm p\to-\bm p\), so the total is
\[
  \ii\int\frac{\dd^3p}{(2\pi)^3}e^{\ii\bm p\cdot(\bm x-\bm y)}
  =
  \ii\delta^{(3)}(\bm x-\bm y).
\]
}

\section{Fock space}

The vacuum is defined by
\[
  a_{\bm p}\ket0=0
  \quad\text{for all }\bm p.
\]
A one-particle state is
\[
  \ket{\bm p}=a_{\bm p}^\dagger\ket0.
\]
Many-particle states are built by repeated creation:
\[
  \ket{\bm p_1,\ldots,\bm p_n}
  =
  a_{\bm p_1}^\dagger\cdots a_{\bm p_n}^\dagger\ket0.
\]
Because the creation operators commute, the particles are bosons:
\[
  a_{\bm p}^\dagger a_{\bm q}^\dagger\ket0
  =
  a_{\bm q}^\dagger a_{\bm p}^\dagger\ket0.
\]
The full space obtained by all such states is Fock space.

\section{The Hamiltonian}

The free Hamiltonian is
\[
  H=\int\dd^3x\,\frac12\left[\pi^2+|\nabla\phi|^2+m^2\phi^2\right].
\]
Substituting the mode expansion and using Fourier delta functions gives
\[
  H=
  \int\frac{\dd^3p}{(2\pi)^3}
  E_{\bm p}
  \left(a_{\bm p}^\dagger a_{\bm p}
  +\frac12(2\pi)^3\delta^{(3)}(0)\right).
\]
The second term is an infinite zero-point energy.  In flat-space scattering calculations it is usually removed by normal ordering:
\[
  :H:
  =
  \int\frac{\dd^3p}{(2\pi)^3}
  E_{\bm p}a_{\bm p}^\dagger a_{\bm p}.
\]
This does not mean the vacuum energy is never physical.  It means that a constant offset in energy does not affect the scattering processes considered here.

\section{Causality and the commutator}

Relativistic causality requires that measurements at spacelike separated points not influence each other.  For a scalar field this is expressed as
\[
  [\phi(x),\phi(y)]=0
  \quad\text{if }(x-y)^2<0.
\]
Using the mode expansion,
\[
  [\phi(x),\phi(y)]
  =
  \int\frac{\dd^3p}{(2\pi)^3}\frac1{2E_{\bm p}}
  \left(e^{-\ii p\cdot(x-y)}-e^{\ii p\cdot(x-y)}\right).
\]
This function is Lorentz invariant.  For spacelike separation, choose a frame where \(x^0-y^0=0\).  Then the integrand is odd under \(\bm p\to-\bm p\), so the integral vanishes.  The argument uses both the relativistic energy \(E_{\bm p}\) and the presence of positive and negative frequency parts in the field.

\section{The Feynman propagator}

The time-ordered two-point function is
\[
  \Delta_F(x-y)=\vev{\order\phi(x)\phi(y)}.
\]
For \(x^0>y^0\),
\[
  \Delta_F(x-y)=
  \int\frac{\dd^3p}{(2\pi)^3}\frac1{2E_{\bm p}}
  e^{-\ii p\cdot(x-y)}.
\]
For \(y^0>x^0\), the sign of \(x-y\) reverses.  Both cases are summarized by
\[
  \Delta_F(x)
  =
  \int\frac{\dd^4p}{(2\pi)^4}
  \frac{\ii e^{-\ii p\cdot x}}{p^2-m^2+\ii\epsilon}.
\]
The small \(+\ii\epsilon\) tells the contour how to pass the poles.  It is the boundary condition that turns a Green function into the Feynman propagator.

\section*{Exercises}
\addcontentsline{toc}{section}{Exercises}

\begin{exercise}
Show that \((\Boxop+m^2)\Delta_F(x)=-\ii\delta^{(4)}(x)\) with the convention above.
\begin{answer}
Act with \(\Boxop+m^2\) on the exponential.  It supplies \((-p^2+m^2)\), leaving \(-\ii\) times the Fourier representation of \(\delta^{(4)}(x)\), up to the harmless \(\epsilon\) limit.
\end{answer}
\end{exercise}

\begin{exercise}
Explain why commuting creation operators imply symmetric two-particle states.
\begin{answer}
Exchanging the labels changes \(a_{\bm p}^\dagger a_{\bm q}^\dagger\ket0\) to \(a_{\bm q}^\dagger a_{\bm p}^\dagger\ket0\), which is equal if the operators commute.
\end{answer}
\end{exercise}

